% Le théorème de Hille-Yosida
% Compile avec pdflatex, xelatex ou lualatex.
\documentclass[11pt,a4paper]{article}

\usepackage{iftex}
\ifPDFTeX
  \usepackage[utf8]{inputenc}
  \usepackage[T1]{fontenc}
  \IfFileExists{lmodern.sty}{\usepackage{lmodern}}{}
  \usepackage[french]{babel}
\else
  \usepackage{fontspec}
  \usepackage{polyglossia}
  \setdefaultlanguage{french}
\fi
\usepackage[margin=2.6cm]{geometry}
\usepackage{amsmath,amssymb,amsthm,mathtools}
\usepackage{booktabs}
\usepackage{enumitem}
\usepackage{microtype}
\usepackage{tikz}
% Ce qu'une `tikzpicture` des fiches peut utiliser, en plus de TikZ.
%
% Lu deux fois : par chaque document (\input après \usepackage{tikz}), pour le
% PDF, et par scripts/tikz.mjs, qui compile chaque figure seule pour la vue web.
% Une bibliothèque ou une couleur ajoutée ici atteint les deux ; ajoutée
% ailleurs, elle n'en atteint qu'un, et le PDF et l'écran ne disent plus la
% même chose.

\usetikzlibrary{calc,arrows.meta,patterns,decorations.pathreplacing}

% La couleur d'accent du site, pour ce qu'une figure veut faire ressortir.
\definecolor{accent}{HTML}{38539D}
\definecolor{accent2}{HTML}{A04B2E}

\usepackage[hidelinks]{hyperref}

% ---------- Environnements ----------
\theoremstyle{plain}
\newtheorem{theoreme}{Théorème}[section]
\newtheorem{proposition}[theoreme]{Proposition}
\newtheorem{lemme}[theoreme]{Lemme}
\newtheorem{corollaire}[theoreme]{Corollaire}
\theoremstyle{definition}
\newtheorem{definition}[theoreme]{Définition}
\newtheorem{exemple}[theoreme]{Exemple}
\newtheorem{exercice}{Exercice}
\theoremstyle{remark}
\newtheorem{remarque}[theoreme]{Remarque}

% ---------- Notations ----------
\newcommand{\R}{\mathbb{R}}
\newcommand{\N}{\mathbb{N}}
\newcommand{\Id}{I}
\newcommand{\dt}{\frac{\mathrm{d}}{\mathrm{d}t}}
\DeclareMathOperator{\Ima}{Im}

\title{Le théorème de Hille-Yosida\\[2pt]
\large Résoudre une équation d'évolution en résolvant des équations stationnaires}
\author{Michel Hua, d'après une note de décembre 2009}
\date{Octobre 2026}

\begin{document}
\maketitle

\begin{abstract}
Le problème de Dirichlet $-\Delta u+\lambda u=f$ se résout par le théorème de Lax-Milgram ; une équation différentielle $u'=F(u)$ à second membre lipschitzien, par celui de Cauchy-Lipschitz. L'équation de la chaleur $\partial_tu-\Delta u=0$ échappe aux deux : elle n'est pas stationnaire, et le laplacien n'est pas lipschitzien, ni même borné. Le théorème de Hille-Yosida la résout en combinant les deux idées. On demande à l'opérateur $A$ d'être monotone et de rendre solubles les équations stationnaires $u+\lambda Au=f$, ce que Lax-Milgram fournit en pratique ; on remplace $A$ par une famille d'opérateurs bornés qui l'approchent, ses régularisées de Yosida, auxquelles Cauchy-Lipschitz s'applique ; et des estimations d'énergie, indépendantes de l'approximation, permettent de passer à la limite. On suit le chapitre VII du livre de Brezis \cite{brezis}, en détaillant toutes les preuves, puis on applique le résultat à l'équation de la chaleur.
\end{abstract}

\paragraph{Leçons concernées.} 222 (Exemples d'équations aux dérivées partielles linéaires), 213 (Espaces de Hilbert. Exemples d'applications), 205 (Espaces complets. Exemples et applications), 208 (Espaces vectoriels normés, applications linéaires continues), 220 (Équations différentielles ordinaires).

\paragraph{Développements possibles.} Le théorème de Hille-Yosida \ref{thm:hy}, avec les estimations du lemme \ref{lem:decroissance} et de la proposition \ref{prop:cauchy} ; la proposition \ref{prop:maxmono}, où le point fixe de Banach propage la surjectivité de $I+\lambda A$ à tout $\lambda>0$ ; la résolution de l'équation de la chaleur (théorème \ref{thm:chaleur}). Le fichier \texttt{HilleYosida.lean} vérifie avec Mathlib les points 1 à 4 de la proposition \ref{prop:maxmono} (le point 4 est pris comme donnée du modèle), les points 1 à 3 de la proposition \ref{prop:yosida}, et le fait que la convergence $J_\lambda v\to v$ pour tout $v\in H$ force la densité de $D(A)$.

\paragraph{Prérequis.} Espaces de Hilbert, théorème de représentation de Riesz, théorème du point fixe de Banach ; pour la dernière section, espaces de Sobolev $H^1_0(\Omega)$ et $H^2(\Omega)$ et formule de Green.

\paragraph{Notations.} $H$ est un espace de Hilbert réel, de produit scalaire $(\cdot,\cdot)$ et de norme $|\cdot|$. Un opérateur linéaire non borné est une application linéaire $A : D(A)\to H$ définie sur un sous-espace $D(A)\subset H$, son \emph{domaine}. On munit $D(A)$ de la norme du graphe $|v|+|Av|$. On note $\|T\|$ la norme d'un opérateur borné $T$ de $H$ dans $H$.

\tableofcontents

% ==========================================================
\section{Deux théorèmes, et une équation qui leur échappe}
% ==========================================================

\paragraph{Lax-Milgram pour le problème de Dirichlet.} Soit $\Omega$ un ouvert de $\R^n$, $f\in L^2(\Omega)$ et $\lambda\geq0$. On cherche $u$ telle que $-\Delta u+\lambda u=f$ dans $\Omega$ et $u=0$ sur $\partial\Omega$. En multipliant par une fonction test et en intégrant par parties, on obtient la formulation variationnelle : trouver $u\in H^1_0(\Omega)$ telle que
\[
  \int_\Omega\nabla u\cdot\nabla v+\lambda\int_\Omega uv=\int_\Omega fv\qquad\text{pour tout }v\in H^1_0(\Omega).
\]

\begin{theoreme}[Lax-Milgram]\label{thm:lm}
Soit $a$ une forme bilinéaire continue sur un espace de Hilbert $V$, \emph{coercive} : il existe $C>0$ tel que $a(v,v)\geq C\|v\|_V^2$ pour tout $v\in V$. Pour toute forme linéaire continue $\ell$ sur $V$, il existe un unique $u\in V$ tel que $a(u,v)=\ell(v)$ pour tout $v\in V$, et $\|u\|_V\leq\frac1C\|\ell\|_{V'}$.
\end{theoreme}

La preuve combine le théorème de représentation de Riesz, qui transforme $a$ et $\ell$ en un opérateur et un vecteur, et le théorème du point fixe de Banach. Avec $V=H^1_0(\Omega)$, la coercivité de $a(u,v)=\int\nabla u\cdot\nabla v+\lambda\int uv$ est immédiate si $\lambda>0$ ; pour $\lambda=0$ et $\Omega$ borné, elle découle de l'inégalité de Poincaré. Le problème de Dirichlet a donc une unique solution faible.

\paragraph{L'équation de la chaleur.} Elle s'écrit
\[
  \frac{\partial u}{\partial t}-\Delta u=0\quad\text{dans }]0,+\infty[\times\Omega,\qquad u=0\ \text{sur }\partial\Omega,\qquad u(0,\cdot)=u_0 .
\]
C'est un problème de Cauchy, avec une donnée initiale, et non plus un problème stationnaire. La formulation variationnelle ne marche plus : le terme $\int\frac{\partial v}{\partial t}v$ n'est pas un carré, et la coercivité de Lax-Milgram disparaît.

On peut aussi voir l'équation comme une équation différentielle ordinaire $u'=\Delta u$, à valeurs dans un espace de fonctions. Le théorème de Cauchy-Lipschitz demande alors que $u\mapsto\Delta u$ soit lipschitzienne, c'est-à-dire, pour un opérateur linéaire, borné. Or le laplacien ne l'est pas : sur $]0,\pi[$, $|\Delta\sin(kx)|=k^2|\sin(kx)|$, et $k$ est arbitraire.

\paragraph{Le cadre abstrait.} On cherche donc à résoudre, pour un opérateur linéaire non borné $A$ de $H$,
\begin{equation}\label{eq:evol}
  \frac{\mathrm{d}u}{\mathrm{d}t}+Au=0\quad\text{sur }[0,+\infty[,\qquad u(0)=u_0 ,
\end{equation}
l'équation de la chaleur correspondant à $H=L^2(\Omega)$ et $A=-\Delta$. En dimension finie, la solution est $u(t)=e^{-tA}u_0$. Mais la série de l'exponentielle n'a aucun sens pour un opérateur non borné, puisqu'elle demande d'appliquer $A^k$ à $u_0$ pour tout $k$.

Le théorème de Hille-Yosida dit quelles hypothèses sur $A$ remplacent la coercivité et le caractère lipschitzien. Elles sont de deux natures : une inégalité, la \emph{monotonie} $(Av,v)\geq0$, qui joue le rôle de la coercivité ; et la possibilité de résoudre les équations stationnaires $u+\lambda Au=f$, qui est exactement ce que Lax-Milgram sait faire. La section 2 étudie ces opérateurs, la section 3 construit leurs approximations bornées, la section 4 démontre le théorème, et la section 5 l'applique à la chaleur.

% ==========================================================
\section{Opérateurs maximaux monotones}
% ==========================================================

\begin{definition}
Un opérateur linéaire $A : D(A)\subset H\to H$ est \emph{monotone} si $(Av,v)\geq0$ pour tout $v\in D(A)$. Il est \emph{maximal monotone} si de plus $\Ima(I+A)=H$, c'est-à-dire si pour tout $f\in H$ l'équation stationnaire $u+Au=f$ a une solution $u\in D(A)$.
\end{definition}

En dimension finie, une matrice $A$ telle que $(Av,v)\geq0$ est automatiquement maximale monotone. En effet $I+A$ est injective, car $v+Av=0$ entraîne $|v|^2+(Av,v)=0$, donc $v=0$ ; et une application linéaire injective d'un espace de dimension finie dans lui-même est surjective. En dimension infinie, la surjectivité est une vraie hypothèse : c'est elle qui contient l'existence de solutions.

L'inégalité $(Av,v)\geq0$ se lit comme une estimation d'énergie. Si $u$ est une solution régulière de \eqref{eq:evol}, alors
\[
  \frac12\dt|u(t)|^2=\Big(\frac{\mathrm{d}u}{\mathrm{d}t},u\Big)=-(Au,u)\leq0 :
\]
la norme de la solution décroît. Pour l'équation de la chaleur, c'est la décroissance de $\int u^2$, une forme de dissipation de l'énergie thermique.

\begin{proposition}[Brezis, proposition VII.1]\label{prop:maxmono}
Soit $A$ un opérateur maximal monotone.
\begin{enumerate}[nosep]
\item $D(A)$ est dense dans $H$.
\item Pour tout $f\in H$, l'équation $u+Au=f$ a une unique solution $u\in D(A)$, et $|u|\leq|f|$.
\item $A$ est fermé : si $u_n\in D(A)$, $u_n\to u$ et $Au_n\to f$, alors $u\in D(A)$ et $Au=f$.
\item Pour tout $\lambda>0$, $\lambda A$ est maximal monotone : $I+\lambda A$ est une bijection de $D(A)$ sur $H$, et son inverse $J_\lambda=(I+\lambda A)^{-1}$ est un opérateur borné de norme $\|J_\lambda\|\leq1$.
\end{enumerate}
\end{proposition}

Dans chaque point, la monotonie sert de la même façon : on prend le produit scalaire de l'équation avec l'inconnue, et le terme $(Av,v)$, positif, se laisse oublier. Seul le point 4 demande une idée de plus, le point fixe de Banach.

\begin{proof}
1. Soit $f\in H$ orthogonal à $D(A)$. Par maximalité, il existe $v_0\in D(A)$ tel que $v_0+Av_0=f$. Alors $0=(f,v_0)=|v_0|^2+(Av_0,v_0)\geq|v_0|^2$, donc $v_0=0$, puis $f=0$.

2. Si $u$ et $\bar u$ sont deux solutions, $w=u-\bar u$ vérifie $w+Aw=0$, d'où $|w|^2+(Aw,w)=0$ et $w=0$. Et $|u|^2+(Au,u)=(f,u)$ donne $|u|^2\leq|f|\,|u|$, soit $|u|\leq|f|$.

3. On a $u_n+Au_n\to u+f$. D'après le point 2, $J_1=(I+A)^{-1}$ est linéaire et de norme au plus $1$, donc continu, et $u_n=J_1(u_n+Au_n)\to J_1(u+f)$. Ainsi $u=J_1(u+f)$ : $u\in D(A)$ et $u+Au=u+f$, c'est-à-dire $Au=f$.

4. La monotonie de $\lambda A$ est claire, et le point 2 s'applique à $\lambda A$ dès que $\lambda A$ est maximal : $J_\lambda$ existe et $\|J_\lambda\|\leq1$. Reste la surjectivité de $I+\lambda A$. Supposons-la acquise pour un $\lambda_0>0$, et montrons-la pour tout $\lambda>\lambda_0/2$. L'équation $u+\lambda Au=f$ équivaut, en multipliant par $\lambda_0/\lambda$, à
\[
  u+\lambda_0Au=\frac{\lambda_0}\lambda f+\Big(1-\frac{\lambda_0}\lambda\Big)u,
  \qquad\text{soit}\qquad
  u=J_{\lambda_0}\Big[\frac{\lambda_0}\lambda f+\Big(1-\frac{\lambda_0}\lambda\Big)u\Big].
\]
Le membre de droite est une application de $H$ dans $H$, lipschitzienne de constante $\big|1-\frac{\lambda_0}\lambda\big|\,\|J_{\lambda_0}\|\leq\big|1-\frac{\lambda_0}\lambda\big|$. Cette constante est $<1$ dès que $\lambda>\lambda_0/2$. Le théorème du point fixe de Banach donne alors un point fixe $u$, qui est dans $D(A)$ puisque $J_{\lambda_0}$ est à valeurs dans $D(A)$, et qui résout l'équation. Partant de $\lambda_0=1$, on obtient la surjectivité pour $\lambda>\frac12$, puis pour $\lambda>\frac14$, et par récurrence pour tout $\lambda>0$.
\end{proof}

\begin{remarque}
Le point 4 est le seul endroit où intervient la complétude de $H$, à travers le point fixe de Banach. C'est aussi lui qui rend la définition commode : il suffit de résoudre une seule équation stationnaire, $u+Au=f$, pour savoir résoudre toutes les équations $u+\lambda Au=f$.
\end{remarque}

\begin{exemple}\label{ex:diag}
Soit $H=\ell^2(\N)$ et $(a_k)$ une suite de réels positifs, non bornée, par exemple $a_k=k^2$. L'opérateur diagonal $A(x_k)=(a_kx_k)$, de domaine $D(A)=\{x\in\ell^2 : \sum a_k^2x_k^2<\infty\}$, est monotone, puisque $(Ax,x)=\sum a_kx_k^2\geq0$. Il est maximal : l'équation $x+Ax=f$ a pour solution $x_k=f_k/(1+a_k)$, qui est bien dans $D(A)$ puisque $a_k|x_k|\leq|f_k|$. Il n'est pas borné, puisque $|Ae_k|=a_k$ pour le $k$-ième vecteur de base. Cet exemple est l'équation de la chaleur sur $]0,\pi[$ vue dans la base des $\sin(kx)$ (section 5).
\end{exemple}

% ==========================================================
\section{La régularisée de Yosida}
% ==========================================================

On voudrait approcher $A$ par des opérateurs bornés, pour pouvoir leur appliquer Cauchy-Lipschitz. La résolvante $J_\lambda=(I+\lambda A)^{-1}$ est bornée, et $\frac1\lambda(I-J_\lambda)$ ressemble formellement à $A$ : si $\lambda A$ était petit, on aurait $J_\lambda\approx I-\lambda A$.

\begin{definition}
Soit $A$ maximal monotone et $\lambda>0$. On appelle \emph{résolvante} de $A$ l'opérateur $J_\lambda=(I+\lambda A)^{-1}$, et \emph{régularisée de Yosida} de $A$ l'opérateur
\[
  A_\lambda=\frac1\lambda(I-J_\lambda) .
\]
\end{definition}

Pour l'opérateur diagonal de l'exemple \ref{ex:diag}, $J_\lambda$ et $A_\lambda$ sont diagonaux, de coefficients $\frac1{1+\lambda a_k}$ et
\[
  \frac1\lambda\Big(1-\frac1{1+\lambda a_k}\Big)=\frac{a_k}{1+\lambda a_k} .
\]
La régularisée garde les petits coefficients presque intacts et écrase les grands sous le plafond $\frac1\lambda$ (figure \ref{fig:yosida}). Quand $\lambda\to0$, le plafond monte et $A_\lambda$ rejoint $A$ coefficient par coefficient.

\begin{figure}[h]
\centering
\begin{tikzpicture}[scale=0.85,>={Stealth[length=2.4mm]}]
  \draw[->,gray] (0,0) -- (6.4,0) node[below] {$a$};
  \draw[->,gray] (0,0) -- (0,4.6);
  \draw[dashed,gray] (0,0) -- (4.4,4.4) node[above right] {$a$};
  \foreach \l/\c in {1/accent, 0.5/accent2, 0.25/black} {
    \draw[\c,thick,domain=0:6.2,samples=80] plot (\x,{\x/(1+\l*\x)});
    \draw[\c,dotted] (0,1/\l) -- (6.2,1/\l);
  }
  \node[accent,right] at (6.2,0.86) {$\lambda=1$};
  \node[accent2,right] at (6.2,1.51) {$\lambda=\frac12$};
  \node[right] at (6.2,2.4) {$\lambda=\frac14$};
  \node[left] at (0,1) {$1$};
  \node[left] at (0,2) {$2$};
  \node[left] at (0,4) {$4$};
\end{tikzpicture}
\caption{Pour un opérateur diagonal de coefficients $a\geq0$, la régularisée de Yosida a pour coefficients $\frac{a}{1+\lambda a}$. Ils restent sous le plafond $\frac1\lambda$ (en pointillé), ce qui rend $A_\lambda$ borné, et tendent vers $a$ quand $\lambda\to0$.}\label{fig:yosida}
\end{figure}

La proposition suivante dit que ce comportement est général.

\begin{proposition}[Brezis, proposition VII.2]\label{prop:yosida}
Soit $A$ maximal monotone et $\lambda>0$.
\begin{enumerate}[nosep]
\item $A_\lambda v=A(J_\lambda v)$ pour tout $v\in H$, et $A_\lambda v=J_\lambda(Av)$ pour tout $v\in D(A)$.
\item $|A_\lambda v|\leq|Av|$ pour tout $v\in D(A)$.
\item $J_\lambda v\to v$ quand $\lambda\to0$, pour tout $v\in H$.
\item $A_\lambda v\to Av$ quand $\lambda\to0$, pour tout $v\in D(A)$.
\item $(A_\lambda v,v)\geq\lambda|A_\lambda v|^2\geq0$ pour tout $v\in H$ : $A_\lambda$ est monotone.
\item $\|A_\lambda\|\leq\frac1\lambda$.
\end{enumerate}
\end{proposition}

\begin{proof}
1. Par définition de $J_\lambda$, on a $v=J_\lambda v+\lambda A(J_\lambda v)$, soit $A(J_\lambda v)=\frac1\lambda(v-J_\lambda v)=A_\lambda v$. Si de plus $v\in D(A)$, on écrit $Av=\frac1\lambda\big[(I+\lambda A)v-v\big]$, et l'on applique $J_\lambda$ : $J_\lambda Av=\frac1\lambda(v-J_\lambda v)=A_\lambda v$.

2. Découle de $A_\lambda v=J_\lambda(Av)$ et de $\|J_\lambda\|\leq1$.

3. Si $v\in D(A)$, alors $|v-J_\lambda v|=\lambda|A_\lambda v|\leq\lambda|Av|\to0$. Dans le cas général, on utilise la densité de $D(A)$ (proposition \ref{prop:maxmono}). Soit $\varepsilon>0$ et $v_1\in D(A)$ tel que $|v-v_1|<\varepsilon$. Alors
\[
  |J_\lambda v-v|\leq|J_\lambda v-J_\lambda v_1|+|J_\lambda v_1-v_1|+|v_1-v|\leq2\varepsilon+|J_\lambda v_1-v_1| ,
\]
donc $\limsup_{\lambda\to0}|J_\lambda v-v|\leq2\varepsilon$ pour tout $\varepsilon>0$.

4. On applique le point 3 à $Av$, puisque $A_\lambda v=J_\lambda(Av)$.

5. On découpe $v=(v-J_\lambda v)+J_\lambda v=\lambda A_\lambda v+J_\lambda v$, d'où, avec le point 1,
\[
  (A_\lambda v,v)=\lambda|A_\lambda v|^2+(A(J_\lambda v),J_\lambda v)\geq\lambda|A_\lambda v|^2 .
\]

6. Par Cauchy-Schwarz et le point 5, $\lambda|A_\lambda v|^2\leq|A_\lambda v|\,|v|$, donc $|A_\lambda v|\leq\frac1\lambda|v|$.
\end{proof}

\begin{remarque}
Le point 3 est le seul qui demande la densité de $D(A)$. La note de 2009 dont cette fiche est issue l'affirmait pour tout $v\in H$ sans cet argument ; le calcul direct n'y suffit pas, car $Av$ n'a de sens que pour $v\in D(A)$. La densité ne peut d'ailleurs pas être évitée : comme $J_\lambda v$ est toujours dans $D(A)$, la convergence $J_\lambda v\to v$ pour tout $v\in H$ entraîne que $D(A)$ est dense. Le point 3 est donc équivalent au point 1 de la proposition \ref{prop:maxmono}, et le fichier Lean le vérifie.
\end{remarque}

Ainsi $(A_\lambda)_{\lambda>0}$ est une famille d'opérateurs bornés et monotones qui converge simplement vers $A$ sur $D(A)$. Leur norme explose quand $\lambda\to0$, comme le plafond de la figure \ref{fig:yosida}, mais pour chaque $\lambda$ fixé le théorème de Cauchy-Lipschitz s'applique.

\begin{theoreme}[Cauchy-Lipschitz]\label{thm:cl}
Soit $E$ un espace de Banach et $F : E\to E$ une application $L$-lipschitzienne. Pour tout $u_0\in E$, il existe une unique $u\in C^1([0,+\infty[;E)$ telle que $u'=F(u)$ et $u(0)=u_0$.
\end{theoreme}

Avec $F=-A_\lambda$, linéaire et de norme au plus $\frac1\lambda$, on obtient pour chaque $\lambda>0$ une unique solution $u_\lambda\in C^1([0,+\infty[;H)$ du problème approché
\begin{equation}\label{eq:approche}
  \frac{\mathrm{d}u_\lambda}{\mathrm{d}t}+A_\lambda u_\lambda=0,\qquad u_\lambda(0)=u_0 .
\end{equation}
Comme $A_\lambda$ est borné, c'est simplement $u_\lambda(t)=e^{-tA_\lambda}u_0$. Il reste à faire tendre $\lambda$ vers $0$, et tout le travail consiste à obtenir des estimations qui ne dépendent pas de $\lambda$.

% ==========================================================
\section{Le théorème de Hille-Yosida}
% ==========================================================

\begin{theoreme}[Hille-Yosida ; Brezis, théorème VII.4]\label{thm:hy}
Soit $A$ un opérateur maximal monotone dans un espace de Hilbert $H$. Pour tout $u_0\in D(A)$, il existe une unique fonction
\[
  u\in C^1([0,+\infty[;H)\cap C([0,+\infty[;D(A))
\]
solution de \eqref{eq:evol}. De plus, pour tout $t\geq0$,
\[
  |u(t)|\leq|u_0|\qquad\text{et}\qquad\Big|\frac{\mathrm{d}u}{\mathrm{d}t}(t)\Big|=|Au(t)|\leq|Au_0| .
\]
\end{theoreme}

\begin{remarque}
La note de 2009 énonçait le théorème pour tout $u_0\in H$. C'est trop demander : si $u\in C([0,+\infty[;D(A))$, alors $u_0=u(0)\in D(A)$, et la borne $|Au(t)|\leq|Au_0|$ n'a de sens que si $u_0\in D(A)$. Pour $u_0\in H$ quelconque, on obtient une solution en un sens affaibli (remarque \ref{rem:semigroupe}).
\end{remarque}

L'intérêt du théorème est de ramener un problème d'évolution à la vérification de deux propriétés de $A$, dont l'une est une équation \emph{stationnaire}. On le démontre en six étapes, comme Brezis.

\paragraph{Étape 1 : unicité.} Si $u$ et $\bar u$ sont deux solutions, $w=u-\bar u$ vérifie $w'+Aw=0$ et $w(0)=0$. Comme $w\in C^1([0,+\infty[;H)$, la fonction $|w|^2$ est dérivable, et
\[
  \frac12\dt|w(t)|^2=(w'(t),w(t))=-(Aw(t),w(t))\leq0 .
\]
Donc $|w(t)|$ décroît ; nulle en $t=0$, elle reste nulle.

L'existence est la partie délicate. On note $u_\lambda$ la solution du problème approché \eqref{eq:approche}, avec $u_0\in D(A)$.

\paragraph{Étape 2 : estimations uniformes.}

\begin{lemme}[Brezis, lemme VII.1]\label{lem:decroissance}
Soit $w\in C^1([0,+\infty[;H)$ telle que $w'+A_\lambda w=0$. Alors les fonctions $t\mapsto|w(t)|$ et $t\mapsto|w'(t)|=|A_\lambda w(t)|$ sont décroissantes sur $[0,+\infty[$.
\end{lemme}

\begin{proof}
On a $\frac12\dt|w|^2=(w',w)=-(A_\lambda w,w)\leq0$ par la monotonie de $A_\lambda$ (proposition \ref{prop:yosida}, point 5). Comme $A_\lambda$ est borné, $w=e^{-tA_\lambda}w(0)$ est de classe $C^\infty$, et $w'$ vérifie la même équation $(w')'+A_\lambda w'=0$ ; ce qui précède s'applique à $w'$.
\end{proof}

Appliqué à $u_\lambda$, avec le point 2 de la proposition \ref{prop:yosida}, ce lemme donne, pour tous $t\geq0$ et $\lambda>0$,
\begin{equation}\label{eq:estim}
  |u_\lambda(t)|\leq|u_0|,\qquad
  \Big|\frac{\mathrm{d}u_\lambda}{\mathrm{d}t}(t)\Big|=|A_\lambda u_\lambda(t)|\leq|A_\lambda u_0|\leq|Au_0| .
\end{equation}
C'est ici que sert l'hypothèse $u_0\in D(A)$ : la vitesse de $u_\lambda$ est contrôlée par $|Au_0|$, indépendamment de $\lambda$.

\paragraph{Étape 3 : la famille $(u_\lambda)$ est de Cauchy.}

\begin{proposition}\label{prop:cauchy}
Pour tous $\lambda,\mu>0$ et $t\geq0$,
\[
  |u_\lambda(t)-u_\mu(t)|\leq\sqrt{2(\lambda+\mu)t}\;|Au_0| .
\]
\end{proposition}

L'idée est d'écrire l'équation vérifiée par $u_\lambda-u_\mu$ et d'estimer sa norme comme à l'étape 1. Le terme $(A_\lambda u_\lambda-A_\mu u_\mu,u_\lambda-u_\mu)$ n'a plus de signe, puisque deux opérateurs différents interviennent ; mais on peut le ramener à $A$, qui est le même pour les deux, à une erreur près de taille $\lambda+\mu$.

\begin{proof}
On soustrait les deux équations et l'on prend le produit scalaire avec $u_\lambda-u_\mu$ :
\[
  \frac12\dt|u_\lambda-u_\mu|^2=-(A_\lambda u_\lambda-A_\mu u_\mu,\,u_\lambda-u_\mu) .
\]
Notons $a=A_\lambda u_\lambda$ et $b=A_\mu u_\mu$. Par définition de $A_\lambda$, $u_\lambda=J_\lambda u_\lambda+\lambda a$, et de même $u_\mu=J_\mu u_\mu+\mu b$. Donc
\[
  (a-b,\,u_\lambda-u_\mu)=(a-b,\,\lambda a-\mu b)+(a-b,\,J_\lambda u_\lambda-J_\mu u_\mu) .
\]
Le second terme est positif : d'après la proposition \ref{prop:yosida}, $a=A(J_\lambda u_\lambda)$ et $b=A(J_\mu u_\mu)$, donc il vaut $\big(A(J_\lambda u_\lambda-J_\mu u_\mu),J_\lambda u_\lambda-J_\mu u_\mu\big)\geq0$ par monotonie de $A$. Le premier vaut $\lambda|a|^2+\mu|b|^2-(\lambda+\mu)(a,b)\geq-(\lambda+\mu)|a|\,|b|$. Avec \eqref{eq:estim}, $|a|,|b|\leq|Au_0|$, d'où
\[
  \frac12\dt|u_\lambda-u_\mu|^2\leq(\lambda+\mu)|Au_0|^2 .
\]
On intègre de $0$ à $t$, en utilisant $u_\lambda(0)=u_\mu(0)=u_0$ : $|u_\lambda(t)-u_\mu(t)|^2\leq2(\lambda+\mu)t\,|Au_0|^2$.
\end{proof}

Brezis obtient la constante $2\sqrt{(\lambda+\mu)t}$, en majorant le premier terme plus grossièrement ; la valeur de la constante ne joue aucun rôle. Pour chaque $t$, $(u_\lambda(t))$ est de Cauchy quand $\lambda\to0$ et converge vers une limite $u(t)$. En faisant tendre $\mu$ vers $0$, on obtient $|u_\lambda(t)-u(t)|\leq\sqrt{2\lambda t}\,|Au_0|$ : la convergence est uniforme sur tout intervalle $[0,T]$, et $u\in C([0,+\infty[;H)$.

\paragraph{Étape 4 : convergence des dérivées, si $u_0\in D(A^2)$.} On suppose de plus $Au_0\in D(A)$. Posons $v_\lambda=u_\lambda'$. Elle vérifie la même équation $v_\lambda'+A_\lambda v_\lambda=0$, avec la donnée initiale $v_\lambda(0)=-A_\lambda u_0$, qui dépend de $\lambda$. Le calcul de l'étape 3 n'utilisait l'égalité des données initiales qu'au moment d'intégrer ; il donne ici
\[
  \frac12\dt|v_\lambda-v_\mu|^2\leq(\lambda+\mu)\,|A_\lambda v_\lambda|\,|A_\mu v_\mu| .
\]
Par le lemme \ref{lem:decroissance}, $|A_\lambda v_\lambda(t)|\leq|A_\lambda v_\lambda(0)|=|A_\lambda A_\lambda u_0|$. Or $A_\lambda u_0=J_\lambda Au_0$, et comme $A_\lambda$ commute avec $J_\lambda$ et $Au_0\in D(A)$, on a $A_\lambda A_\lambda u_0=J_\lambda A_\lambda Au_0=J_\lambda J_\lambda A^2u_0$, de norme au plus $|A^2u_0|$. En intégrant,
\[
  |v_\lambda(t)-v_\mu(t)|^2\leq|A_\lambda u_0-A_\mu u_0|^2+2(\lambda+\mu)t\,|A^2u_0|^2 ,
\]
et le premier terme tend vers $0$ (proposition \ref{prop:yosida}, point 4). Les dérivées $u_\lambda'$ convergent donc uniformément sur tout $[0,T]$.

\paragraph{Étape 5 : existence, si $u_0\in D(A^2)$.} Par les étapes 3 et 4, $u_\lambda\to u$ et $u_\lambda'$ converge, uniformément sur tout $[0,T]$ ; donc $u\in C^1([0,+\infty[;H)$ et $u_\lambda'\to u'$. On écrit le problème approché sous la forme $u_\lambda'(t)+A(J_\lambda u_\lambda(t))=0$. Or $J_\lambda u_\lambda(t)\to u(t)$, car
\[
  |J_\lambda u_\lambda(t)-u(t)|\leq|J_\lambda u_\lambda(t)-J_\lambda u(t)|+|J_\lambda u(t)-u(t)|\leq|u_\lambda(t)-u(t)|+|J_\lambda u(t)-u(t)|\longrightarrow0 ,
\]
et $A(J_\lambda u_\lambda(t))=-u_\lambda'(t)\to-u'(t)$. Comme $A$ est fermé, $u(t)\in D(A)$ et $Au(t)=-u'(t)$. Enfin $t\mapsto Au(t)=-u'(t)$ est continue, donc $u\in C([0,+\infty[;D(A))$. Les estimations \eqref{eq:estim} passent à la limite et donnent $|u(t)|\leq|u_0|$ et $|u'(t)|\leq|Au_0|$.

\paragraph{Étape 6 : existence pour $u_0\in D(A)$.} On approche $u_0$ par des données de $D(A^2)$.

\begin{lemme}[Brezis, lemme VII.2]\label{lem:densite}
$D(A^2)$ est dense dans $D(A)$ pour la norme du graphe.
\end{lemme}

\begin{proof}
Soit $u_0\in D(A)$ et $\lambda>0$. Le vecteur $\tilde u_0=J_\lambda u_0$ est dans $D(A)$, et $A\tilde u_0=A_\lambda u_0=J_\lambda Au_0$ est encore dans $D(A)$ : donc $\tilde u_0\in D(A^2)$. Quand $\lambda\to0$, $\tilde u_0=J_\lambda u_0\to u_0$ et $A\tilde u_0=J_\lambda Au_0\to Au_0$, par le point 3 de la proposition \ref{prop:yosida}.
\end{proof}

Soit $(u_{0n})$ une suite de $D(A^2)$ telle que $u_{0n}\to u_0$ et $Au_{0n}\to Au_0$, et $u_n$ la solution donnée par l'étape 5. Pour deux indices $m$ et $n$, $u_n-u_m$ est la solution de donnée initiale $u_{0n}-u_{0m}\in D(A^2)$ ; les estimations de l'étape 5 donnent, pour tout $t\geq0$,
\[
  |u_n(t)-u_m(t)|\leq|u_{0n}-u_{0m}|,\qquad|u_n'(t)-u_m'(t)|\leq|Au_{0n}-Au_{0m}| .
\]
Les suites $(u_n)$ et $(u_n')$ sont donc de Cauchy uniformément sur $[0,+\infty[$ : $u_n\to u$ et $u_n'\to u'$ uniformément, avec $u\in C^1([0,+\infty[;H)$. Comme $Au_n=-u_n'\to-u'$ et $A$ est fermé, $u(t)\in D(A)$ et $u'+Au=0$ ; les estimations passent à la limite. Cela termine la preuve du théorème \ref{thm:hy}. \qed

\begin{remarque}[semi-groupe de contractions]\label{rem:semigroupe}
Pour $t\geq0$, notons $S(t)u_0=u(t)$ la solution issue de $u_0\in D(A)$. L'application $S(t)$ est linéaire, par unicité, et $|S(t)u_0|\leq|u_0|$. Comme $D(A)$ est dense, elle se prolonge en un opérateur continu de $H$ dans $H$, de norme au plus $1$. La famille $(S(t))_{t\geq0}$ vérifie $S(0)=I$, $S(t+s)=S(t)S(s)$, et $S(t)u_0\to u_0$ quand $t\to0$, pour tout $u_0\in H$ : c'est un \emph{semi-groupe continu de contractions}, l'analogue de $e^{-tA}$. Pour $u_0\in H\setminus D(A)$, $t\mapsto S(t)u_0$ est continue mais n'est en général dérivable en aucun point : c'est une solution « généralisée » de \eqref{eq:evol}. Réciproquement, tout semi-groupe continu de contractions provient d'un unique opérateur maximal monotone ; c'est la forme complète du théorème de Hille-Yosida (Brezis, remarque VII.5).
\end{remarque}

\begin{remarque}[effet régularisant]\label{rem:regularisant}
Si $A$ est de plus autoadjoint, la situation est bien meilleure : pour tout $u_0\in H$, la solution généralisée vérifie
\[
  u\in C([0,+\infty[;H)\cap C^1(]0,+\infty[;H)\cap C(]0,+\infty[;D(A)) ,
\]
et $|Au(t)|\leq\frac1t|u_0|$ pour $t>0$ (Brezis, théorème VII.7, que nous admettons). La solution est régulière dès que $t>0$, même si la donnée ne l'est pas. Le laplacien de Dirichlet est autoadjoint, et c'est l'effet régularisant bien connu de l'équation de la chaleur.
\end{remarque}

% ==========================================================
\section{Application à l'équation de la chaleur}
% ==========================================================

Soit $\Omega$ un ouvert borné de $\R^n$, de frontière de classe $C^2$. On prend $H=L^2(\Omega)$, et
\[
  Au=-\Delta u,\qquad D(A)=H^2(\Omega)\cap H^1_0(\Omega) .
\]
La condition $u\in H^1_0(\Omega)$ contient la condition au bord $u=0$ sur $\partial\Omega$.

\begin{proposition}\label{prop:laplacien}
L'opérateur $A=-\Delta$ de domaine $H^2(\Omega)\cap H^1_0(\Omega)$ est maximal monotone dans $L^2(\Omega)$.
\end{proposition}

\begin{proof}
\emph{Monotonie.} Pour $v\in D(A)$, la formule de Green donne $(Av,v)=-\int_\Omega v\Delta v=\int_\Omega|\nabla v|^2\geq0$, le terme de bord disparaissant puisque $v=0$ sur $\partial\Omega$.

\emph{Maximalité.} Soit $f\in L^2(\Omega)$. On résout $u-\Delta u=f$ en deux temps. D'abord, par le théorème de Lax-Milgram dans $V=H^1_0(\Omega)$ avec $a(u,v)=\int\nabla u\cdot\nabla v+\int uv$, qui est coercive de constante $1$ pour la norme de $H^1$, il existe un unique $u\in H^1_0(\Omega)$ solution faible. Ensuite, le théorème de régularité elliptique (Brezis, théorème IX.25, que nous admettons) montre que $u\in H^2(\Omega)$, et l'équation $u-\Delta u=f$ a alors lieu presque partout. Donc $u\in D(A)$ et $u+Au=f$.
\end{proof}

On voit sur cette preuve le rôle de chaque théorème. La maximalité de $A$ est exactement le problème de Dirichlet de la section 1, avec $\lambda=1$, et c'est Lax-Milgram qui la fournit ; la coercivité ne demande même pas l'inégalité de Poincaré, puisque le terme $\int uv$ est présent.

\begin{theoreme}\label{thm:chaleur}
Pour tout $u_0\in H^2(\Omega)\cap H^1_0(\Omega)$, il existe une unique fonction
\[
  u\in C^1\big([0,+\infty[;L^2(\Omega)\big)\cap C\big([0,+\infty[;H^2(\Omega)\cap H^1_0(\Omega)\big)
\]
telle que $\partial_tu-\Delta u=0$ et $u(0)=u_0$. De plus $\|u(t)\|_{L^2}\leq\|u_0\|_{L^2}$ et $\|\Delta u(t)\|_{L^2}\leq\|\Delta u_0\|_{L^2}$ pour tout $t\geq0$.
\end{theoreme}

\begin{proof}
On applique le théorème \ref{thm:hy} à l'opérateur de la proposition \ref{prop:laplacien}.
\end{proof}

Comme $A$ est autoadjoint, la remarque \ref{rem:regularisant} s'applique : pour une donnée initiale seulement dans $L^2(\Omega)$, la solution généralisée est dans $H^2(\Omega)\cap H^1_0(\Omega)$ dès que $t>0$.

\begin{exemple}\label{ex:intervalle}
Sur $\Omega=]0,\pi[$, les fonctions $e_k(x)=\sqrt{2/\pi}\,\sin(kx)$, $k\geq1$, forment une base hilbertienne de $L^2(]0,\pi[)$, et $-e_k''=k^2e_k$. Dans cette base, $A$ est l'opérateur diagonal de l'exemple \ref{ex:diag} avec $a_k=k^2$. Si $u_0=\sum_kb_ke_k$, la solution est
\[
  u(t,x)=\sum_{k\geq1}b_k\,e^{-k^2t}\,e_k(x) .
\]
Chaque mode décroît d'autant plus vite que sa fréquence est grande : les oscillations s'effacent d'abord (figure \ref{fig:chaleur}). La régularisée de Yosida remplace $e^{-k^2t}$ par $e^{-tk^2/(1+\lambda k^2)}$ : elle amortit à peine les hautes fréquences, ce qui est le prix de la bornitude.
\end{exemple}

\begin{figure}[h]
\centering
\begin{tikzpicture}[xscale=2,yscale=1.9]
  \draw[->,gray] (0,0) -- (3.4,0) node[below] {$x$};
  \draw[->,gray] (0,-0.2) -- (0,1.6);
  \node[below] at (3.1416,0) {$\pi$};
  \node[left] at (0,1) {$1$};
  \draw[gray!50] (-0.04,1) -- (0.04,1);
  \foreach \t/\c/\y/\lab in {0/black/1.45/{t=0}, 0.05/accent2/1.2/{t=0{,}05}, 0.2/accent/0.95/{t=0{,}2}, 0.6/gray/0.7/{t=0{,}6}} {
    \draw[\c,thick,domain=0:3.1416,samples=90] plot (\x,{exp(-\t)*sin(\x r)+0.5*exp(-9*\t)*sin(3*\x r)});
    \draw[\c,thick] (3.55,\y) -- (3.8,\y) node[right,black] {$\lab$};
  }
\end{tikzpicture}
\caption{Solution de l'équation de la chaleur sur $]0,\pi[$ pour $u_0(x)=\sin x+\frac12\sin3x$ : $u(t,x)=e^{-t}\sin x+\frac12e^{-9t}\sin3x$. Le mode $\sin3x$ décroît neuf fois plus vite que $\sin x$, et le profil s'arrondit avant de s'éteindre.}\label{fig:chaleur}
\end{figure}

\begin{remarque}[l'équation de la chaleur rétrograde]
Le signe de $A$ n'est pas un détail. Pour $\partial_tu+\Delta u=0$, c'est-à-dire $A=\Delta$, qui n'est pas monotone, la donnée $u_0=\frac1k\sin(kx)$ donne la solution $u(t,x)=\frac1ke^{k^2t}\sin(kx)$. Quand $k\to\infty$, la donnée tend vers $0$ dans $L^2$, mais pour tout $t>0$ fixé, $\|u(t)\|_{L^2}\to\infty$ : la solution ne dépend pas continûment de la donnée. Le problème est mal posé au sens d'Hadamard, et aucun théorème du type Hille-Yosida ne peut le résoudre.
\end{remarque}

% ==========================================================
\section{Bilan}
% ==========================================================

L'équation de la chaleur échappait à Lax-Milgram, faute d'être stationnaire, et à Cauchy-Lipschitz, faute de second membre borné. Le théorème de Hille-Yosida prend à chacun ce dont il a besoin.

\begin{center}
\renewcommand{\arraystretch}{1.25}
\begin{tabular}{@{}p{3.2cm}p{3.6cm}p{3.6cm}p{4.2cm}@{}}
\toprule
\textbf{Théorème} & \textbf{Problème} & \textbf{Hypothèse} & \textbf{Outil de preuve} \\
\midrule
Lax-Milgram & stationnaire, $a(u,v)=\ell(v)$ & coercivité & Riesz et point fixe de Banach \\
Cauchy-Lipschitz & évolution, $u'=F(u)$ & $F$ lipschitzienne & point fixe de Banach \\
Hille-Yosida & évolution, $u'+Au=0$ & $A$ monotone et $I+A$ surjectif & régularisée de Yosida, estimations d'énergie, $A$ fermé \\
\bottomrule
\end{tabular}
\end{center}

La monotonie joue le rôle de la coercivité : c'est elle qui donne les estimations d'énergie, de l'unicité à la proposition \ref{prop:cauchy}. La maximalité est une équation stationnaire, résolue en pratique par Lax-Milgram. Et la régularisée de Yosida permet d'utiliser Cauchy-Lipschitz à $\lambda$ fixé, avant de passer à la limite.

% ==========================================================
\section{Exercices}
% ==========================================================

\begin{exercice}\label{exo:dimfinie}
Soit $A\in M_n(\R)$ telle que $(Av,v)\geq0$ pour tout $v\in\R^n$. Montrer que $\|e^{-tA}\|\leq1$ pour tout $t\geq0$, pour la norme d'opérateur associée à la norme euclidienne. Calculer $J_\lambda$ et $A_\lambda$ pour $A=\begin{pmatrix}0&1\\-1&0\end{pmatrix}$, et vérifier que $A_\lambda\to A$.
\end{exercice}

\begin{exercice}\label{exo:transport}
Soit $H=L^2(\R)$ et $Au=u'$, de domaine $D(A)=H^1(\R)$. Montrer que $A$ est maximal monotone, en résolvant explicitement $u+u'=f$. Vérifier que la solution de \eqref{eq:evol} est $u(t,x)=u_0(x-t)$, et que $|u(t)|=|u_0|$ pour tout $t$.
\end{exercice}

\begin{exercice}\label{exo:yosida}
Dans l'exemple \ref{ex:intervalle}, on note $u_\lambda$ la solution du problème approché \eqref{eq:approche}. Montrer que $|u_\lambda(t)-u(t)|\to0$ pour tout $t\geq0$ et tout $u_0\in L^2(]0,\pi[)$, en travaillant mode par mode.
\end{exercice}

\begin{exercice}\label{exo:decalage}
Soit $A$ maximal monotone et $c\in\R$. Résoudre $u'+Au+cu=0$, $u(0)=u_0\in D(A)$, en posant $v(t)=e^{ct}u(t)$.
\end{exercice}

% ==========================================================
\section{Corrigés}
% ==========================================================

\paragraph{Exercice 1.} Soit $u(t)=e^{-tA}u_0$. Alors $\frac12\dt|u|^2=-(Au,u)\leq0$, donc $|e^{-tA}u_0|\leq|u_0|$ pour tout $u_0$. Pour $A=\begin{pmatrix}0&1\\-1&0\end{pmatrix}$, on a $(Av,v)=0$ et $A^2=-I$. Donc $(I+\lambda A)(I-\lambda A)=I-\lambda^2A^2=(1+\lambda^2)I$, et
\[
  J_\lambda=\frac{I-\lambda A}{1+\lambda^2},\qquad
  A_\lambda=\frac1\lambda(I-J_\lambda)=\frac{\lambda I+A}{1+\lambda^2}\xrightarrow[\lambda\to0]{}A .
\]
Ici $A$ est antisymétrique, et $e^{-tA}$ est une rotation : la norme est conservée, pas seulement contrôlée.

\paragraph{Exercice 2.} Pour $u\in H^1(\R)$, $(Au,u)=\int_\R uu'=\frac12\big[u^2\big]_{-\infty}^{+\infty}=0$, car une fonction de $H^1(\R)$ tend vers $0$ à l'infini. Pour $f\in L^2(\R)$, la fonction
\[
  u(x)=\int_{-\infty}^xe^{-(x-y)}f(y)\,\mathrm{d}y
\]
est la convolution de $f$ avec $g(x)=e^{-x}\mathbf{1}_{x>0}$, qui est intégrable. Par l'inégalité de Young, $u\in L^2(\R)$ et $|u|\leq|f|$ ; puis $u'=f-u\in L^2(\R)$, donc $u\in H^1(\R)$ et $u+u'=f$. Enfin $u(t,x)=u_0(x-t)$ vérifie $\partial_tu+\partial_xu=0$, et la translation conserve la norme $L^2$. Comme à l'exercice 1, $(Au,u)=0$ donne une évolution qui conserve la norme.

\paragraph{Exercice 3.} Dans la base $(e_k)$, $A_\lambda$ est diagonal de coefficients $\frac{k^2}{1+\lambda k^2}$, donc
\[
  u_\lambda(t)=\sum_kb_k\,e^{-tk^2/(1+\lambda k^2)}\,e_k ,
\]
et
\[
  |u_\lambda(t)-u(t)|^2=\sum_kb_k^2\Big(e^{-tk^2/(1+\lambda k^2)}-e^{-tk^2}\Big)^2 .
\]
Chaque terme tend vers $0$ quand $\lambda\to0$ et est majoré par $b_k^2$, terme général d'une série convergente. Par convergence dominée pour les séries, la somme tend vers $0$. On n'a pas eu besoin de $u_0\in D(A)$ ; en revanche, la convergence n'est pas uniforme en $u_0$, puisque pour $u_0=e_k$ l'écart à $t=1$ vaut $e^{-k^2/(1+\lambda k^2)}-e^{-k^2}$, proche de $e^{-1/\lambda}$ quand $k$ est grand.

\paragraph{Exercice 4.} Si $u$ est une solution, $v(t)=e^{ct}u(t)$ vérifie $v'=e^{ct}(u'+cu)=-e^{ct}Au=-Av$ et $v(0)=u_0$. Réciproquement, le théorème \ref{thm:hy} donne une unique $v$, et $u(t)=e^{-ct}v(t)$ est l'unique solution, avec $|u(t)|\leq e^{-ct}|u_0|$. Le cas $c<0$ autorise une croissance exponentielle contrôlée ; c'est le cadre des opérateurs « $\omega$-monotones », $(Av,v)\geq-\omega|v|^2$.

\begin{thebibliography}{9}
\bibitem{brezis} H. Brezis, \emph{Analyse fonctionnelle. Théorie et applications}, Masson, 1983, chapitre VII (pp. 101-110) pour le théorème de Hille-Yosida, chapitre X pour l'équation de la chaleur. Version anglaise augmentée : \emph{Functional Analysis, Sobolev Spaces and Partial Differential Equations}, Springer, 2011, chapitre 7.
\bibitem{ciarlet} P. Ciarlet, É. Lunéville, \emph{La méthode des éléments finis : de la théorie à la pratique. I. Concepts généraux}, Les Presses de l'ENSTA. Le théorème de Lax-Milgram et le problème de Dirichlet.
\bibitem{pazy} A. Pazy, \emph{Semigroups of Linear Operators and Applications to Partial Differential Equations}, Springer, 1983. La théorie des semi-groupes dans un espace de Banach, et la forme générale du théorème de Hille-Yosida.
\end{thebibliography}

\paragraph{Source.} Note « Le théorème d'Hille-Yosida », écrite en décembre 2009 d'après le chapitre VII du livre de Brezis. La note s'arrêtait avant la preuve de l'existence ; cette fiche complète la preuve (étapes 2 à 6), corrige l'énoncé (donnée initiale dans $D(A)$ et non dans $H$) et ajoute l'application à l'équation de la chaleur.

\end{document}
